%% kapa: exam=1 round=28 year=2017 period=1 form=A subject=economics subject_ko=경제학원론 seq=29 number=69
%% source: Q-Net 감정평가사 기출문제 (cst00309) artlSeq=5199603 file=extracted/1008/5199603/2017년 제28회 감정평가사 1차시험 문제지/1교시 A형.pdf
%% conversion: kordoc 4.24.51 -> tools/gen1.py (auto); review=proofread-sonnet; flags=none
\begin{question}{69}
인플레이션에 관한 설명으로 옳지 않은 것은?
\begin{choices}
  \item 프리드만(M. Friedman)에 따르면 인플레이션은 언제나 화폐적 현상이다.
  \item 정부가 화폐공급을 통해 얻게 되는 추가적인 재정수입이 토빈세(Tobin tax)이다.
  \item 비용상승 인플레이션은 총수요관리를 통한 단기 경기안정화정책을 어렵게 만든다.
  \item 예상하지 못한 인플레이션은 채권자에서 채무자에게로 소득재분배를 야기한다.
  \item 인플레이션이 예상되는 경우에도 메뉴비용(menu cost)이 발생할 수 있다.
\end{choices}
\end{question}
\answer{2}
